% ==============================================================================
% capitulos/01_introducao.tex
% Capítulo 1 — Introdução
% ==============================================================================

\chapter{Introdução}
\label{cap:introducao}

\section{Contextualização}
\label{sec:contextualizacao}

A transformação digital dos ambientes fabris, impulsionada pelo paradigma da Indústria 4.0, pela Internet das Coisas (IoT) e pela sua vertente industrial (IIoT), tem aproximado de forma crescente os sistemas de Tecnologia de Operação (OT) das plataformas de Tecnologia da Informação (IT). Historicamente, Controladores Lógicos Programáveis (CLPs) operavam confinados em redes de campo isoladas, dedicados exclusivamente ao controle sequencial, ao intertravamento de máquinas e à supervisão local \cite{moraes_castrucci}. Nas arquiteturas conectadas modernas, entretanto, espera-se que esses controladores atuem como nós inteligentes de borda (\textit{edge devices}), expondo variáveis operacionais em tempo quase real para pipelines analíticas capazes de identificar anomalias, sintetizar diagnósticos contextualizados e retroalimentar ações supervisórias ao processo \cite{chandola_2009, stouffer_nist_ot}.

Controladores de nova geração como o WEG PLC500 --- programável em ambiente CODESYS conforme a norma IEC 61131-3 e equipado com duas interfaces Gigabit Ethernet independentes com suporte nativo ao protocolo Modbus TCP/IP --- oferecem uma ponte robusta entre o chão de fábrica e serviços de software modernos \cite{weg_plc500_manual, iec_61131_3}. Paralelamente, ferramentas de programação visual orientadas a fluxos de eventos (\textit{flow-based programming}), com destaque para o Node-RED, consolidaram-se em soluções de IoT e automação pela facilidade de integração com protocolos industriais por meio do pacote \texttt{node-red-contrib-modbus}, conectividade com bancos de dados relacionais e capacidade de orquestrar serviços analíticos e modelos de Inteligência Artificial (IA) \cite{nodered_docs, nodered_contrib_modbus}.

\section{Histórico do Projeto, Redefinição de Escopo e Enquadramento na Situação de Aprendizagem}
\label{sec:historico_ntp}

Este projeto desenvolve-se no âmbito de uma Situação de Aprendizagem (SA) de Internet das Coisas (IoT) de escopo livre, cuja diretriz pedagógica permite a proposição de arquiteturas e problemas reais de engenharia desde que submetidos à aprovação prévia do professor responsável. Esse formato metodológico visa estimular a autonomia projetual, a capacidade de diagnóstico de viabilidade técnica e a integração ponta a ponta entre hardware, protocolo de comunicação em rede, processamento de dados e atuação.

Dentro dessa modalidade de escopo livre, a proposta inicial submetida ao docente tinha um tema distinto no domínio de redes e sistemas embarcados conectados: o desenvolvimento de um servidor de tempo \textit{Network Time Protocol} (NTP) Stratum~1 utilizando uma antena de recepção \textit{Global Navigation Satellite System} (GNSS) acoplada a um microcontrolador ESP32 \cite{rfc5905_ntp, espressif_esp32}. Na concepção original, o receptor GNSS forneceria ao ESP32 tanto as sentenças de navegação no padrão NMEA (via interface serial UART) quanto o pulso elétrico de sincronismo de alta precisão \textit{Pulse Per Second} (PPS) por interrupção de hardware, permitindo ao microcontrolador disciplinar seu relógio local e atender requisições NTPv4 na rede local.

Durante a fase inicial de avaliação experimental e prototipação em laboratório, contudo, identificaram-se restrições práticas severas que inviabilizaram a continuidade e a validação confiável do servidor NTP em ambiente acadêmico:
\begin{enumerate}
  \item Inviabilidade de recepção de sinal GNSS em ambiente interno (\textit{indoor}): os sinais de radiofrequência transmitidos pelas constelações GNSS na banda~L chegam à superfície terrestre com potência extremamente baixa (tipicamente inferior a $-125\,\text{dBm}$). No interior do laboratório, a atenuação provocada pelas lajes de concreto armado e pela estrutura metálica da edificação impedia a aquisição de efemérides e o travamento (\textit{3D fix}) de satélites suficientes para habilitar o sinal PPS, enquanto medições próximas a janelas sofriam forte degradação por multicaminho (\textit{multipath}).
  \item Limitações de infraestrutura física e de hardware: a instalação permanente de uma antena GNSS ativa externa no telhado do prédio, interligada ao laboratório por cabo coaxial de baixa perda com proteção contra surtos, mostrou-se impraticável dentro das restrições físicas e logísticas do ambiente, somando-se a dificuldades de disponibilidade e estabilidade dos módulos periféricos do protótipo com ESP32.
\end{enumerate}

Diante desse diagnóstico técnico, realizou-se uma pivotagem completa de tema, redirecionando o projeto para um sistema de Internet Industrial das Coisas (IIoT) e análise de dados integralmente reproduzível e validável em bancada \textit{indoor}, sem dependência de visada externa para satélites. Em estrita observância às regras da Situação de Aprendizagem de IoT de escopo livre, o novo escopo --- fundamentado no controlador industrial WEG PLC500 integrado via Modbus TCP/IP à plataforma Node-RED, com botoeiras simulando sensores de campo e indicadores LED simulando atuadores --- foi formalmente apresentado e contou com a aprovação prévia do professor, o que garante plena conformidade do projeto com o estabelecido para a avaliação da disciplina. Ademais, um princípio técnico relevante da proposta original foi preservado na nova arquitetura: o rigor com a consistência temporal das amostras (\textit{timestamping}) e com a avaliação analítica de eventos em janelas de tempo bem definidas.

\section{Definição do Problema e Motivação}
\label{sec:problema}

Na arquitetura adotada após a redefinição de escopo, a integração entre uma pipeline analítica com recursos de IA no Node-RED e um CLP industrial impõe dois desafios centrais de engenharia:

\begin{enumerate}
  \item Fidelidade de simulação segura em bancada: antes do acoplamento a motores de potência, válvulas pressurizadas ou processos térmicos reais, é indispensável dispor de uma bancada didático-industrial controlada que represente os sinais discretos de campo por meio de botoeiras e LEDs em 24\,Vcc, permitindo exercitar trepidação mecânica de contatos, condição de nível crítico, sobretemperatura, parada de emergência e desconexão de rede sem risco físico.
  \item Convivência entre análise inteligente não determinística e segurança determinística: fluxos executados na camada IT --- incluindo \textit{polling} sobre TCP/IP, consultas a banco de dados e chamadas a Modelos de Linguagem Grandes (LLMs) --- apresentam latência variável e comportamento probabilístico. Por conseguinte, toda decisão analítica deve ser validada por contratos de dados estritos (\textit{JSON Schema}) e permanecer subordinada a uma camada determinística de intertravamento \textit{fail-safe} executada localmente no WEG PLC500 \cite{stouffer_nist_ot}.
\end{enumerate}

A motivação deste trabalho reside em documentar uma solução de engenharia completa que enderece ambos os problemas, combinando a confiabilidade determinística do WEG PLC500, a padronização do protocolo Modbus TCP/IP e a capacidade de orquestração e análise híbrida do Node-RED.

\section{Objetivos}
\label{sec:objetivos}

\subsection{Objetivo Geral}
Especificar e implementar a arquitetura técnica, os contratos de comunicação, a lógica determinística de controle, a pipeline analítica híbrida (estatística, aprendizado de máquina via FastAPI/scikit-learn e diagnóstico cognitivo via LLM) e a verificação automatizada \textit{Software-in-the-Loop} (SIL) de um sistema de supervisão de sensores simulados por botoeiras em um controlador WEG PLC500 integrado via Modbus TCP/IP à plataforma Node-RED.

\subsection{Objetivos Específicos}
\begin{itemize}
  \item Definir e documentar um cenário didático-industrial estruturado sobre as entradas e saídas digitais nativas de 24\,Vcc PNP do WEG PLC500 (quatro botoeiras simulando sensores de processo e quatro LEDs com resistores de $2{,}2\,\text{k}\Omega$ simulando atuadores e sinalizadores), contemplando lógica normalmente fechada (NF) para parada de emergência e artefatos completos de engenharia de hardware (\texttt{hardware/bom.csv}, \texttt{hardware/wiring\_table.md}, diagramas CircuiTikZ/TikZ e esquemático KiCad~8 \texttt{plc500\_bench.kicad\_sch}).
  \item Especificar a topologia de rede OT/IT segmentada e o mapa completo de 23 registradores Modbus TCP/IP em endereçamento base zero (\textit{Discrete Inputs}, \textit{Coils}, \textit{Input Registers} e \textit{Holding Registers}), incluindo mecanismo de \textit{heartbeat} e \textit{watchdog} de comunicação.
  \item Implementar a lógica determinística modular no WEG PLC500 em \textit{Structured Text} (IEC 61131-3) no CODESYS V3.5 (\texttt{GVL\_Modbus.st}, \texttt{FB\_DebounceCounter.st}, \texttt{FB\_WatchdogHeartbeat.st}, \texttt{FB\_SafetyInterlock.st} e \texttt{PLC\_PRG.st}, exportados também em formato PLCopen TC6 XML v2.01), abrangendo filtragem de \textit{debounce}, contadores modulares de 16~bits, máquina de estados finitos e intertravamento local com precedência absoluta sobre comandos remotos.
  \item Desenvolver a pipeline analítica híbrida em cinco abas na plataforma Node-RED v4 com o pacote \texttt{node-red-contrib-modbus}, combinando persistência em PostgreSQL~16, extração de atributos em janelas deslizantes de $60\,\text{s}$ (\textit{stride} de $15\,\text{s}$), microsserviço de inferência em aprendizado de máquina via FastAPI e scikit-learn (\texttt{IsolationForest} e \texttt{HistGradientBoostingClassifier}), diagnóstico contextual via agente LLM governado por \textit{JSON Schema} Draft 2020-12 e emissão automática de relatórios operacionais.
  \item Estabelecer e verificar o plano formal de validação (\texttt{CT-01} a \texttt{CT-10}) utilizando um gêmeo digital Modbus TCP determinístico com ciclo de varredura de $10\,\text{ms}$ (\texttt{plc\_simulator.py}), critérios quantitativos de aceitação e matriz de análise de riscos técnicos.
\end{itemize}

\section{Escopo, Delimitações e Conformidade com a Situação de Aprendizagem}
\label{sec:escopo}

Por tratar-se de uma Situação de Aprendizagem de IoT de escopo livre com aprovação prévia do professor, a conformidade do projeto com o estabelecido é assegurada simultaneamente no plano formal (validação docente prévia da redefinição de escopo após os ensaios preliminares de bancada) e no plano técnico-pedagógico, uma vez que a solução contempla de ponta a ponta todas as camadas estruturantes de um ecossistema IoT/IIoT:
\begin{itemize}
  \item Sensoriamento, atuação física e controle de borda (\textit{Edge/OT}): aquisição digital em 24\,Vcc PNP com filtragem de \textit{debounce} ($50\,\text{ms}$), contagem modular de eventos, máquina de estados e intertravamento local \textit{fail-safe} no WEG PLC500, acompanhada de documentação elétrica executiva em KiCad~8 e CircuiTikZ;
  \item Conectividade e interoperabilidade em rede: comunicação cliente/servidor sobre Ethernet segmentada (\texttt{ETH2}) utilizando protocolo aberto Modbus TCP/IP na porta~502 com supervisão contínua de enlace (\textit{heartbeat}/\textit{watchdog});
  \item Orquestração, persistência e processamento analítico híbrido (\textit{IT/IoT Platform}): fluxos orientados a eventos em cinco abas no Node-RED (\texttt{node-red-\allowbreak contrib-\allowbreak modbus}), armazenamento histórico e forense em PostgreSQL~16 (quatro tabelas e três visões operacionais), detecção estatística e inferência supervisionada/não supervisionada via FastAPI (\texttt{Isolation\allowbreak Forest} com ROC-AUC de teste de $0{,}8993$ e \texttt{HistGradient\allowbreak Boosting\allowbreak Classifier} com acurácia de $99{,}73\%$ e Macro-F1 de $0{,}9956$ sobre conjunto sintético de $18\,\text{h}$ a $2\,\text{Hz}$ totalizando $129.600$ amostras brutas e $3.688$ janelas deslizantes com lacunas de purga $\ge 60\,\text{s}$) e diagnóstico assistido por LLM governado por JSON Schema;
  \item Retroalimentação em malha fechada, gêmeo digital e observabilidade: comando supervisionado dos atuadores (LEDs) com auditoria da máscara real das saídas físicas (\texttt{wMascara\_\allowbreak DO\_\allowbreak Real}), dashboard operacional (\texttt{@flowfuse/\allowbreak node-red-\allowbreak dashboard}) e validação automatizada \textit{Software-in-the-Loop} (\texttt{CT-01} a \texttt{CT-10}) contra o simulador determinístico \texttt{plc\_\allowbreak simulator.py}.
\end{itemize}

Este documento consolida a arquitetura, a especificação técnica, a implementação completa no monorrepo e o protocolo de validação verificado em ambiente \textit{Software-in-the-Loop} e preparado para comissionamento físico em bancada.

A camada física de campo restringe-se a sinais discretos (entradas e saídas digitais nativas em 24\,Vcc PNP), sem utilização de módulos analógicos neste ciclo, mantendo todavia a arquitetura pronta para expansão futura com cartões da família MOD/RUW100 do WEG PLC500 \cite{weg_plc500_manual}.

\section{Organização do Documento}
\label{sec:organizacao}

O restante do documento estrutura-se em seis capítulos e um apêndice:
\begin{itemize}
  \item O Capítulo~\ref{cap:fundamentacao} aborda a fundamentação técnica sobre o WEG PLC500, a norma IEC 61131-3, o protocolo Modbus TCP/IP, a plataforma Node-RED com \texttt{node-red-contrib-modbus}, a análise híbrida de anomalias (estatística, aprendizado de máquina e LLM) e os princípios de segurança OT/IT.
  \item O Capítulo~\ref{cap:arquitetura} detalha a arquitetura em camadas, a topologia de rede, os artefatos de engenharia de hardware (BOM, tabela de fiação, diagramas CircuiTikZ e KiCad~8), o mapeamento das botoeiras e LEDs, o mapa de 23 registradores Modbus TCP/IP e o diagrama de sequência da comunicação.
  \item O Capítulo~\ref{cap:implementacao_clp} especifica a configuração no CODESYS V3.5, a arquitetura modular de blocos funcionais em \textit{Structured Text}, a exportação PLCopen TC6 XML v2.01, a máquina de estados finitos e o código embarcado no WEG PLC500.
  \item O Capítulo~\ref{cap:implementacao_nodered} descreve a implementação de cinco abas no Node-RED v4, a formulação matemática e extração de 10 atributos em janelas deslizantes de $60\,\text{s}$, o microsserviço FastAPI de inferência scikit-learn, a governança do agente LLM com JSON Schema Draft 2020-12 e o esquema relacional com visões operacionais no PostgreSQL~16.
  \item O Capítulo~\ref{cap:plano_validacao} estabelece o protocolo de validação com dez casos de teste (\texttt{CT-01} a \texttt{CT-10}), a verificação automatizada \textit{Software-in-the-Loop} com o gêmeo digital \texttt{plc\_simulator.py}, os critérios quantitativos de aceitação e a análise de riscos técnicos.
  \item O Capítulo~\ref{cap:consideracoes} sintetiza as conclusões de engenharia e apresenta o roteiro de trabalhos futuros.
  \item O Apêndice~\ref{ap:checklist} disponibiliza um \textit{checklist} prático de comissionamento físico e implantação da pilha de software e validação da bancada.
\end{itemize}
